\section{The Virtue Annexed to Justice}

The virtue that pays the unpayable debt has carried the name \emph{religio} since Roman times, and the ancients already disagreed about where the word came from. Cicero derived it from \latin{relegere}, to read over again---the religious man ponders repeatedly what concerns the worship of God. Augustine liked \latin{reeligere}, to choose again what negligence had lost, and elsewhere \latin{religare}, to bind: ``May religion bind us to the one Almighty God.'' Aquinas, reviewing the candidates, declines to arbitrate, because all three etymologies land in the same place: however the word was built, ``it denotes properly a relation to God. For it is He to Whom we ought to be bound as to our unfailing principle; to Whom also our choice should be resolutely directed as to our last end.''\endnote{\emph{Summa theologiae} II-II, q.~81, a.~1, corpus, reviewing Cicero's \latin{relegere} (via Isidore), Augustine's \latin{reeligere} (\emph{De civitate Dei} X.3), and \latin{religare} (\emph{De vera religione} 55); the quoted conclusion is Aquinas's own. Checked 2026-07-24, English at New Advent, Latin at Corpus Thomisticum. ``Unfailing principle'' and ``last end'' deliberately echo the two directions of Part I's argument: the source from which being is presently received, and the good toward which the receiver is ordered.} Re-reading, re-choosing, re-binding: each is a \emph{return}---the trace, inside a Latin noun, of the psalmist's \latin{retribuam}.

That religion is a virtue at all follows from the analysis of the previous section in one step. A virtue is a habit that makes its possessor good and his work good; and ``to render anyone his due has the aspect of good, since by rendering a person his due, one becomes suitably proportioned to him, through being ordered to him in a becoming manner.''\endnote{\emph{Summa theologiae} II-II, q.~81, a.~2, corpus; checked 2026-07-24 as above. The interior argument is worth pausing on: rendering the due is good not because the creditor profits but because the renderer becomes \emph{rightly ordered}---proportioned to reality. Applied to God, this is the whole efficacy of worship, as section 4 develops.} Note where the goodness lands: not on the creditor, who in this case can receive nothing, but on the debtor, who in rendering becomes \emph{rightly ordered}---proportioned to what is actually the case. And that religion is a \emph{special} virtue, distinct from generic justice and from every reverence paid to creatures, follows from the uniqueness of its term: honor is due wherever there is excellence, and ``to God a singular excellence is competent, since He infinitely surpasses all things and exceeds them in every way. Wherefore to Him is special honor due.''\endnote{\emph{Summa theologiae} II-II, q.~81, a.~4, corpus; checked 2026-07-24 as above. The same question grounds the traditional distinction between \latin{latria}, the worship owed to God alone, and the reverence (\latin{dulia}) paid to excellent creatures---a distinction Part I touched at II-II, q.~84, a.~1, and which Augustine's \emph{City of God} X develops at length (section~5 below).} One excellence without peer; therefore one honor without analogue; therefore a virtue with, so to speak, God's name on it---the only moral virtue defined entirely by \emph{whom} it faces.

Then comes the placement that surprises everyone on first meeting, and it is the structural heart of this essay. Religion is \emph{not} a theological virtue. Faith, hope, and charity have God himself as their object---by them the soul reaches out to God, believes him, hopes in him, loves him. Religion does not touch God. Its matter is creatures: words, gestures, hours, vows, animals once, bread and wine, the body's positions---acts and things \emph{offered} to God but never arriving in him the way faith's assent arrives. ``God is related to religion not as matter or object, but as end: and consequently religion is not a theological virtue whose object is the last end, but a moral virtue which is properly about things referred to the end.''\endnote{\emph{Summa theologiae} II-II, q.~81, a.~5, corpus; checked 2026-07-24 as above. The reply to the first objection completes the picture from above: the theological virtues, whose acts do bear on God as proper object, ``by their command\ldots{} cause the act of religion, which performs certain deeds directed to God''---faith, hope, and charity command what religion executes. II-II, q.~81, a.~6 then ranks religion first among the moral virtues as ``approaching nearer to God than the other moral virtues, in so far as its actions are directly and immediately ordered to the honor of God.''} A moral virtue: acquired the way virtues are acquired, exercised in ordinary matter, seated in the will, ranked---first, indeed, but ranked---among the cardinal virtues' companions, below the three that touch God.

The placement does two kinds of work, and both matter for a creature trying to live Part I's conclusion. First, it makes worship \emph{sober}. If adoration belonged to the theological order, one might suspect it of being an effusion---something performed out of the warmth of faith when the warmth is flowing, and honestly omitted when it is not. Annexed to justice, it is instead a matter of what is \emph{owed}: due on the days one feels the vertigo and on the days one feels nothing, the way a debt is owed in every mood. The tradition, asked whether worship is really a species of enthusiasm, answers: no---of debt. Second, and oppositely, the placement keeps worship \emph{humble} in the technical sense. A moral virtue about ``things referred to the end'' cannot pretend that its performances reach, please, improve, or affect God. Religion's acts are arrows shot toward a target they cannot wound or even touch; their whole reality as acts terminates inside creation---in changed words, changed hours, a changed worshipper. Both errors about worship die together here: the error that treats it as optional sentiment, and the error that treats it as effective magic. It is neither. It is justice, doing the one thing justice can do when equality is impossible: rendering the acknowledgment that is due, with a constant and perpetual will.

One more feature of the architecture, easy to miss and impossible to overstate. Religion has, Aquinas says, two ranges of acts. Some it \emph{elicits}---produces directly as its own: adoration, sacrifice, prayer, the acts whose immediate sense is Godward. But others it \emph{commands}: acts produced by other virtues and bent by religion toward the honor of God. His example is deliberate: ``to visit the fatherless and widows in their tribulation is an act of religion as commanding, and an act of mercy as eliciting.''\endnote{\emph{Summa theologiae} II-II, q.~81, a.~1, ad~1; checked 2026-07-24 as above. Aquinas's example quotes James 1:27 (``Religion clean and undefiled before God and the Father, is this: to visit the fatherless and widows in their tribulation: and to keep one's self unspotted from this world,'' Douay--Rheims), which he reads through the eliciting/commanding distinction: mercy performs the visit, religion refers it to God. The essay returns to this joint when Augustine identifies mercy as ``the true sacrifice'' in section~5.} The scriptural sentence behind the example---religion clean and undefiled is to visit the fatherless and widows---has been quoted in every century against ritualism, as if it abolished the altar in favor of the soup kitchen. Aquinas reads it more carefully: mercy performs the visit, religion refers the visit to God, and the same concrete afternoon is both. The debt owed to the source of all good is not paid only in a sanctuary. It annexes the whole moral life, the way a spring pressurizes every pipe in the house. What remains to be seen is what the elicited acts---the Godward ones proper---actually are, and why a creature made of body as well as mind cannot pay its acknowledgment in thought alone.
